% GENERATED by scripts/make_supplementary_tables.py from git 7c55fbc (verbatim text; references resolved, floats in the v5 supplementary material left out: tab:c2_classify, tab:c2_custody, tab:c2_partitions, tab:c2_targets, tab:main)
\section{Earlier Evaluation on the Held-Out Period}
\label{app:earlier}

This appendix reports the evaluation that preceded the confirmation experiment.
It used a pass/fail criterion with thresholds that were later abandoned
(Section~\ref{M-sec:setup:earlier}); the results are reported as they were obtained.
The criterion was replaced because the result rested on one evaluation seed and a
narrow margin in the rarest group, the held-out period had been reused for
further analyses (Section~\ref{app:heldout} of the supplementary material), and
the fixed thresholds were conventions.
A reanalysis of the same period in the format of Section~\ref{M-sec:setup:metrics}
is given at the end of this appendix.

\subsection{Periods, Selection, and Criteria}
\label{app:earlier:protocol}

\begin{table}[!htbp]
\centering
\caption{Periods of the earlier evaluation (month/day, all dates inclusive)}
\footnotesize
\label{tab:earlier_periods}
\begin{tabular}{@{}llll@{}}
\toprule
Period & Training & Calib. & Evaluation segments \\
\midrule
Expl.~1 & 11/04--12/01 & 12/02--12/08 & 12/09--12/22 \\
(2013--14) & & & 12/23--01/05 (holidays) \\
Expl.~2 & 02/03--03/02 & 03/03--03/09 & 03/10--03/23 \\
(2014) & & & 03/24--04/13 \\
 & & & 04/14--04/27 (Easter) \\
Held-out & 05/05--06/01 & 06/02--06/08 & A: 06/09--06/20 \\
(2014) & & & B: 06/21--06/30 \\
\bottomrule
\end{tabular}
\end{table}

We use three disjoint periods (Table~\ref{tab:earlier_periods}), each split into four weeks
of training, one week of calibration, and two or three evaluation segments. The two
exploratory periods were used to develop the methods, to choose the groups, and to
compare candidates; the layout attributes branching and isolation were added after
analyses of the misses on these periods. The third period was held out.

\paragraph*{Method selection and criteria}
Among the candidates (MVP on the uniform and the logarithmic scale,
\texttt{qr\_groups}, and \texttt{qr\_ogd}), we retained those for which, in all
five exploratory segments, no group had a 95\% Wilson upper confidence bound below
\groupBar and whose mean cost was at most \costFactor{} times that of the hindsight
oracle, and chose the cheapest. This selected
\texttt{mvp\_log}, the only candidate meeting both conditions; \texttt{qr\_ogd}
became the secondary method by the tie-break rule among the candidates meeting the
cost condition. The same two conditions are the evaluation criteria on the
held-out period. The value \groupBar{} was in the plan from its first draft,
initially as a bound on the point estimate of the worst group; its present form,
which counts groups significantly below \groupBar, and the cost factor
\costFactor, i.e., at most 20\% more search than the oracle, were set after
exploratory results with an earlier, smaller collection of groups and before the
final groups were evaluated. With the final groups, the exploratory cost of the
selected method was \explMvpLogCostRatioMean{} times that of the oracle. The
selection rule and the criteria were recorded in version control before the
held-out period was evaluated.

\paragraph*{Main results and additional analyses}
The held-out period was evaluated once to obtain the main results
(Appendix~\ref{app:earlier:main}). All further analyses, listed in
Section~\ref{app:heldout} of the supplementary material, were performed after the main results were known;
they reuse the held-out period and, where stated, the exploratory periods, and
were not used to revise the methods or criteria.\footnote{The evaluation code
initially tabulated a slightly different collection of groups (a first-hop
indicator added, the group of all hops omitted). Recomputing all statistics on
the \numGroups{} groups from the stored hop-level results changed no decision
reported here.}

\paragraph*{Criteria of the earlier evaluation}

For each segment and each group with at least \minGroupHops{} hops, we report the
recall with its 95\% Wilson interval~\cite{wilson1927}. We call \groupBar{} the
\emph{bar}. The group criterion requires that no group be significantly below the
bar in either held-out segment; the cost criterion requires a mean cost of at most
\costFactor{} times that of the oracle. The bar lies \barGapPointsWord{} point below the target
because the online methods aim at recall \emph{equal} to the target, so large
groups slightly below it would otherwise be flagged; omitting a multiple-testing
correction makes the criterion stricter. Because not rejecting is not evidence of
adequate recall, we also report the worst group, the count of groups significantly
below \targetRecall, and the number of evaluation seeds for which the criterion
holds. Cost is the mean camera-seconds per hop, relative to the oracle of the same
predictor and segment.


\subsection{Marginal Calibration Misses Hard Groups}
\label{app:earlier:problem}

On held-out segment~B (Fig.~\ref{fig:group_recall} of the supplementary material), split conformal calibration
nearly meets the target on average (\heldStaticRecallB) but leaves
\heldStaticBelowBarB{} groups below the bar, in identifiable situations: the first
hop after a lost track is re-acquired (recall \heldStaticWorstB), hops leaving
sites with high branching (\heldStaticBranchHighRecallB) or high isolation
(\heldStaticIsoHighRecallB), the first hop of a trip
(\heldStaticInitiatedRecallB), and slow exits (\heldStaticSlowRecallB). Hops
leaving sites with low branching or isolation reach about
\heldStaticEasyRecallApproxB: the average is met by over-searching easy hops.
Segment~A shows the same pattern. The failure is not a matter of tuning: the
hindsight oracle, whose single threshold meets the target exactly on the
evaluation segment, still leaves \heldOracleBelowBarA{} and \heldOracleBelowBarB{}
groups below the bar (Table~\ref{S-tab:main}). Which attributes matter depends on
the layout: on a synthetic layout of the same city, exit speed rather than
isolation dominated (Table~\ref{tab:layout}), which suggests targeting many
attributes at once rather than one known difficulty.

\subsection{Main Results on the Held-Out Period}
\label{app:earlier:main}

\texttt{mvp\_log}, selected as primary method under the earlier plan, meets both criteria
(Table~\ref{S-tab:main}): no group is below the bar in either segment (worst groups
\heldMvpLogWorstA{} and \heldMvpLogWorstB), and its cost is \heldMvpLogCostRatioA{}
and \heldMvpLogCostRatioB{} times that of the oracle, both below \costFactor. With a bar on the point
estimate of the worst group, the form of the first draft of the plan, it would not
pass in segment~B. \texttt{qr\_ogd}, selected as secondary method, has a similar cost
(\heldQrOgdCostRatioMean) but leaves one group below the bar in segment~B. The two
are not statistically distinguishable: their worst groups in segment~B,
re-acquisition hops in both cases, have overlapping intervals, and a McNemar test
on the paired outcomes does not reject equality ($p=\heldWorstBPairedP$).

No baseline meets both criteria. The marginal rules (\texttt{static}, ACI,
\texttt{mvp\_global}) and the oracle leave \heldMarginalBelowBarMin{} to
\heldMarginalBelowBarMax{} groups below the bar, Mondrian
\heldMondrianBelowBarA{} and \heldMondrianBelowBarB, and \texttt{qr\_groups} three
in segment~B. MVP on the uniform scale and the CityTrac-style rule meet the group
criterion only at higher cost: \heldMvpGroupsCostRatioA{} and
\heldMvpGroupsCostRatioB{} times the oracle for the former, and
\heldCityTracNinetyFiveCostRatioA{} and \heldCityTracNinetyFiveCostRatioB{} times
for the latter (target \citytracTargetLow); the latter trusts uncalibrated
probabilities and reaches a mean recall of \heldCityTracNinetyFiveRecallApprox.
With the bar at the target itself, every method except the CityTrac-style rule
with target \citytracTargetHigh{} (at \heldCityTracNinetyNineCostRatioA{} and
\heldCityTracNinetyNineCostRatioB{} times the oracle cost) leaves groups
significantly below \targetRecall. For \texttt{mvp\_log}, these shortfalls are at
most \heldMvpLogMaxShortfallPpA{} points in segment~A and, apart from the
re-acquisition group, \heldMvpLogMaxShortfallPpExReacqB{} points in segment~B, as
expected for a method that targets recall equal to the target.

The oracle meets the target only on average. Against two references that meet it
in every group in hindsight (Table~\ref{tab:group_oracle}; additional analysis),
the cost of \texttt{mvp\_log} equals that of a group quantile regression fitted on
the evaluation segment itself (ratios \goMvpLogOverQrA{} and \goMvpLogOverQrB),
while the rule that gives each hop the lowest hindsight threshold of its groups
costs \goGroupMinCostRatioA{} and \goGroupMinCostRatioB{} times the oracle. Online calibration thus reaches the
cost of a group-conditional fit in hindsight, but not its worst-group recall
(\goQrWorst).

A comparison at equal cost is given in Section~\ref{app:earlier:frontier} of the supplementary material.

\subsection{Reanalysis in the Confirmation Format}
\label{app:earlier:reanalysis}

After the confirmation results were known, the held-out period was reanalysed
with the targets, estimates, comparators, hindsight references, and decision
rules of the confirmation experiment, as the locked plan specified in advance.
This is an additional analysis of a period whose results were already known. It
is reported in Section~\ref{supp:reanalysis} of the supplementary material as a
check of consistency between periods; the verdicts of the confirmation period
stand.

\begin{table*}[t]
\centering
\caption{Parameter values of the experiments.}
\label{tab:parameters}
\footnotesize
% AUTO-GENERATED by scripts/make_tables.py -- DO NOT EDIT. Data: generated/data/tab_parameters.csv
\begin{tabular}{@{}p{3.6cm}p{1.9cm}p{3.4cm}p{7.5cm}@{}}
\toprule
Item & Symbol & Value & Note \\
\midrule
\multicolumn{4}{@{}l}{\emph{Data and preprocessing}} \\
Bounding box (lon, lat) &  & [-8.7, -8.55] $\times$ [41.1, 41.2] & a trip is kept if at least one GPS point lies inside \\
Minimum points per trip &  & 3 & trips with fewer points are dropped \\
Maximum speed &  & 50 m/s & any step faster: whole trip dropped \\
Missing-data trips &  & dropped & trips flagged in the data set as having missing points \\
GPS interval / interpolation step &  & 15 s / 1 s & linear interpolation, no map matching \\
\addlinespace[3pt]
\multicolumn{4}{@{}l}{\emph{Cameras}} \\
Cameras in the file &  & 239 & camera layer of the Porto open-data portal (version retrieved in July 2025) \\
Tunnel cameras &  & excluded & cameras flagged as inside a tunnel; cameras without the flag are treated as outside \\
Merging &  & complete linkage, 50 m & 176 sites \\
Observation radius &  & 50 m & disk around each site; at most 4 sites queried per point \\
\addlinespace[3pt]
\multicolumn{4}{@{}l}{\emph{Hops and groups}} \\
Delay bin / horizon / bins & $\Delta$, $T_{\max}$, $B$ & 5 s, 600 s, 120 & no re-appearance within $T_{\max}$: exit \\
Time-of-day edges &  & 0, 6, 10, 16, 20, 24 h & local time of the exit \\
Exit zone &  & 4 sectors & E, N, W, S sector of the last point in the source disk relative to the site, boundaries at $\pm45^\circ$ \\
Exit speed &  & last 4 s in the disk & straight-line distance between the last interpolated point in the source site's disk and the point 4 s earlier (or the first point of the sighting), divided by the elapsed time \\
Branching &  & entropy & entropy of the source site's next-site distribution over training non-exit hops \\
Isolation &  & distance & distance from the source site to the nearest other site \\
Quartile edges &  & training period & speed, branching and isolation, over the training hops; set for each period from its own training period (held-out: 05/05--06/01, 2014, dates inclusive) \\
Source unseen in training &  &  & branching: highest quartile; predictor: flat matrix at $\epsilon_0$ \\
\addlinespace[3pt]
\multicolumn{4}{@{}l}{\emph{Predictor}} \\
Candidate radius & $R_{\mathrm{cand}}$ & 2500 m & plus every next site observed in training \\
Shrinkage & $\lambda$ & 20 & toward the parent key \\
Base probability & $\epsilon_0$ & 0.0001 & prior of the coarsest level \\
Delay smoothing & $\sigma_g$ & 1 bin & Gaussian along the delay axis \\
Context keys &  & source site, exit zone, time of day, exit-speed quartile & backed off to the deepest key seen in training \\
\addlinespace[3pt]
\multicolumn{4}{@{}l}{\emph{Feedback}} \\
Feedback fraction & $\rho$ & 0.3 & random over evaluation hops (evaluation seed) \\
Resolution time &  & arrival / $t_{\mathrm{exit}}+T_{\max}$ & hit resolved at re-appearance, miss at the end of the horizon \\
\addlinespace[3pt]
\multicolumn{4}{@{}l}{\emph{Calibration methods}} \\
ACI & $\gamma_A$, $W$ & 0.0005, 2000 & window initialized with the last 2000 calibration scores; $\alpha_t$ clipped to [-0.5, 1] \\
Mondrian &  & time of day $\times$ first hop of a track $\times$ exit-speed quartile & cells with fewer than 200 calibration hops use the global threshold \\
qr\_groups &  & level $\alpha$, log score & intercept + reference-coded indicators of the six attributes; subsample of at most 80{,}000 calibration hops; score floor 1e-08 \\
qr\_ogd & $\gamma$, clip & 0.005, $\pm$3 & offsets start at 0 (threshold of qr\_groups), log scale \\
MVP (mvp\_log) & $m$, $r$, $\epsilon$ & 40, 1000, 1 & $\eta$ = 0.0319 (25 groups), 0.1167 (1 group) \\
 & $a_{\min}$ & 0.002 & warm start: all calibration hops in time order with immediate feedback \\
\addlinespace[3pt]
\multicolumn{4}{@{}l}{\emph{Evaluation}} \\
Target / group bar & $1-\alpha$ & 0.95 / 0.94 & groups significantly below the bar: 95\% Wilson upper bound \\
Cost criterion &  & 1.20 & mean of A and B, cost / hindsight oracle \\
Minimum group size &  & 200 & smaller groups not judged \\
Evaluation seeds &  & 0 (main); 0, 1, 2, 3, 4 & additional seed runs \\
\bottomrule
\end{tabular}

\end{table*}



\section{Confirmation Experiment: Further Results}
\label{app:confirm2}

This appendix reports the results of the confirmation experiment that the main
text summarizes. All values are from both blocks pooled and averaged over seeds
unless stated otherwise; intervals are the conservative 95\% intervals of the
taxi and day cluster bootstraps.

\paragraph*{Other targets}
Table~\ref{S-tab:c2_targets} extends Table~\ref{M-tab:c2_main} to targets 0.90,
0.93, 0.97, and 0.98.


\paragraph*{All pre-specified comparisons}
Table~\ref{S-tab:c2_classify} classifies every pre-specified comparison of an
online method with a comparator for the three estimates (re-acquisition
shortfall, isolation~Q4 shortfall, and $\bar{D}$) at the five targets, with 99\%
simultaneous intervals, and gives for each pair the number of pairs of target
points at which one method dominated the other: its 95\% upper bound of cost was
below the other's 95\% lower bound, and its 95\% lower bound of worst-group
recall above the other's 95\% upper bound (Sec.~10.5 of the analysis plan in
the supplementary material).
% v4 S10 PENDING (JD): plan attached as supplementary material.


\paragraph*{Blocks}
Fig.~\ref{fig:c2_blocks} repeats Figs.~\ref{M-fig:c2_calibration}
and~\ref{M-fig:c2_cost} for each block, and Table~\ref{tab:c2_blocks} gives the
between-block differences of $\bar{D}$, with 95\% intervals, and of the
re-acquisition shortfall, as point estimates.

\begin{table*}[!htbp]
\centering
\caption{Difference between the two blocks (autumn minus August) in $\bar{D}$ (pp, with 95\% intervals) and in the re-acquisition shortfall (pp, point estimates), per method and target.}
\label{tab:c2_blocks}
\footnotesize
\IfFileExists{generated/tables/tab_c2_blocks.tex}{% source: results/porto_cctv_confirm2/analysis/block_comparison.csv @ 1a806b3; results/porto_cctv_confirm2/analysis/deficits.csv @ 1a806b3
% Dbar: -(diff) of block_comparison.csv (= D_block2 - D_block1, autumn minus August) with [-hi, -lo] (conservative 95%: taxi paired, day independent); re-acquisition: deficit(block2) - deficit(block1) of deficits.csv cell hopkind=2, point only (no interval for this difference in the result files). Block 1, block 2: D_block1, D_block2 of block_comparison.csv (points). Rows: mvp_log, qr_ogd, qr_groups_refit, mondrian_hour_hop_speed (shortened from the full table; Dbar only).
\setlength{\tabcolsep}{3.4pt}
\begin{tabular}{@{}llrrr@{}}
\toprule
Method & Target & Block~1 \Dbar{} & Block~2 \Dbar{} & Difference [95\%] \\
\midrule
\texttt{mvp\_log} & 0.90 & 1.6 & 0.7 & $-$0.8 [$-$2.7, 0.9] \\
 & 0.93 & 2.2 & 0.8 & $-$1.4 [$-$2.9, 0.4] \\
 & 0.95 & 1.7 & 0.7 & $-$1.0 [$-$2.1, 0.6] \\
 & 0.97 & 1.4 & 0.7 & $-$0.6 [$-$1.6, 0.7] \\
 & 0.98 & 1.1 & 0.7 & $-$0.4 [$-$1.3, 0.6] \\
\addlinespace
\texttt{qr\_ogd} & 0.90 & 0.7 & 0.8 & 0.1 [$-$1.6, 1.7] \\
 & 0.93 & 1.0 & 0.8 & $-$0.2 [$-$1.5, 1.5] \\
 & 0.95 & 1.0 & 0.6 & $-$0.4 [$-$1.3, 1.0] \\
 & 0.97 & 0.6 & 0.5 & $-$0.1 [$-$0.7, 0.4] \\
 & 0.98 & 0.5 & 0.6 & 0.2 [$-$0.3, 0.7] \\
\addlinespace
\texttt{qr\_groups\_refit} & 0.90 & 0.5 & 0.5 & 0.0 [$-$1.7, 0.6] \\
 & 0.93 & 0.9 & 0.6 & $-$0.3 [$-$1.9, 0.5] \\
 & 0.95 & 0.7 & 0.4 & $-$0.3 [$-$1.2, 0.4] \\
 & 0.97 & 0.4 & 0.3 & $-$0.1 [$-$0.4, 0.4] \\
 & 0.98 & 0.3 & 0.3 & 0.1 [$-$0.2, 0.7] \\
\addlinespace
Mondrian & 0.90 & 8.2 & 6.8 & $-$1.4 [$-$1.8, $-$0.8] \\
 & 0.93 & 5.3 & 4.2 & $-$1.1 [$-$2.4, $-$0.6] \\
 & 0.95 & 3.9 & 3.0 & $-$0.9 [$-$2.5, 0.6] \\
 & 0.97 & 3.1 & 2.1 & $-$1.0 [$-$2.3, 0.7] \\
 & 0.98 & 2.3 & 1.5 & $-$0.8 [$-$1.9, 0.4] \\
\bottomrule
\end{tabular}
}{%
\fbox{\parbox{0.9\textwidth}{\centering tab\_c2\_blocks --- table from Alex}}}
\end{table*}

\paragraph*{Partitions not used for calibration}
Table~\ref{S-tab:c2_partitions} gives, for the pre-specified alternative partitions,
the paired difference between their largest shortfall and that on the
\numGroups{} groups, and the shortfall on the intersections I1--I6.


\paragraph*{Diagnostics, tuning history, and trial run}
The diagnostics of the confirmation runs (seed ranges, the ACI variants, the
share of hops at $a_{\min}$, and reachable upper bounds) and the tuning history
of each method with the trial run before the lock are reported in
Tables~\ref{tab:c2_diagnostics} and~\ref{tab:c2_tuning} of the supplementary
material.

% ---------------------------------------------------------------------
\subsection{Track Custody (Additional Analysis)}
\label{app:confirm2:custody}

This analysis was decided after the confirmation results were known and was
not part of the locked plan. Table~\ref{S-tab:c2_custody} reports the share of
chains whose first $k$ hops were all found, $k=1,\dots,5$, with the definition of
Table~\ref{tab:custody} (cumulative minimum of the open-loop hits).


\section{Implementation Details and Availability}
\label{app:implementation}

% v4-5 (request 01, item 9): notation table moved from Section III.
\begin{table}[!htbp]
\centering
\caption{Main notation}
\label{tab:notation}
\begin{tabular}{@{}ll@{}}
\toprule
Symbol & Meaning \\
\midrule
$\mathcal{C}$, $N$ & set and number of camera sites \\
$\mathcal{R}_c$ & observation region of site $c$ \\
$t^{\mathrm{exit}}$, $s$ & end time and source site of an episode \\
$x$ & exit context observed at $t^{\mathrm{exit}}$ \\
$\Delta$, $T_{\max}$, $\mathcal{B}$ & bin width, horizon, and bin index set \\
$\mathcal{V}$ & visibility set of a hop \\
$\mathcal{S}$ & search set (camera--time pairs) \\
$\hat{p}(c,b\mid x)$ & estimated visibility probability \\
$\tau$, $\sigma$ & threshold and score of a hop \\
$1-\alpha$ & target recall \\
$\mathcal{G}$ & collection of (overlapping) groups \\
$\rho$ & fraction of hops that provide feedback \\
$p_{\mathrm{det}}$ & per-visit detection probability \\
\midrule
\multicolumn{2}{@{}l}{\emph{Method (Section~\ref{M-sec:method})}} \\
$\mathcal{G}(x)$ & groups that contain context $x$ \\
$q$, $\tau(q)$ & calibration index and its threshold \\
$a_{\min}$ & smallest calibration miss level (log scale) \\
$m$, $B_i$ & number of buckets and $i$-th bucket \\
$V_{G,i}$, $n_{G,i}$ & coverage surplus and count, group $G$, bucket $i$ \\
$\eta$ & MVP learning rate \\
$\theta_G$, $\gamma$ & online group offset and its step size \\
\bottomrule
\end{tabular}
\end{table}

\paragraph*{Data}
The trajectory file contains \numTripsFile{} trips of \numTaxisFile{} distinct taxi
identifiers (the dataset description mentions \numTaxisDescribed). We keep trips with
at least one GPS point inside the bounding box and at least \minTripPointsWord{}
points, and discard trips flagged as having missing data and trips with any step
implying a speed above \maxSpeed~m/s. The trajectories are available from the UCI
Machine Learning Repository under the Creative Commons Attribution 4.0
(CC~BY~4.0) license~\cite{porto2015data}. The camera layer is the dataset
``Sistema de Gest\~ao de Tr\'afego -- CCTV'' of the open-data portal of the
Municipality of Porto~\cite{porto2025cctv} (last updated February~7, 2025);
because the portal was not reachable from outside Portugal, we retrieved the file
through the Internet Archive (snapshot of July~9, 2025). It lists
\numCamerasRaw{} cameras with a tunnel flag; we exclude the
\numTunnelCamerasWord{} camera flagged as inside a tunnel, where GPS positions are
unreliable, and treat cameras without the flag as outside tunnels. The current
version of the layer, released under the Creative Commons CC0~1.0 Universal
public-domain dedication, lists \portalCamerasCurrent{} cameras; we kept the
version used from the start of the study. The code, the configuration files, and
the result summaries will be made available in a public repository upon
publication.
% TODO(JD): repository URL and archive DOI (e.g., Zenodo) at publication.

\paragraph*{Visibility predictor}
The predictor conditions on the keys $k_1(x)=(s)$, $k_2(x)=(s,z)$,
$k_3(x)=(s,z,h)$, and $k_4(x)=(s,z,h,v)$, where $z$ is the exit zone, $h$ the
time-of-day period, and $v$ the exit-speed quartile. For level $d$ and key $k$, let $n_d(k)$ be the number of non-exit training hops
with $k_d(x)=k$, and let
\begin{equation}
  M_d(k)[c,b] = \sum_{t:\,k_d(x_t)=k} \mathbb{1}\bigl[(c,b)\in\mathcal{V}_t\bigr]
  \label{eq:counts}
\end{equation}
for $c\in\mathcal{C}_s$ and $b\in\mathcal{B}$ be the number of those hops in which
the object was visible in site $c$ during bin $b$. To account for small timing
errors, each row of $M_d(k)$ is smoothed along the time axis with a Gaussian
kernel of width $\sigma_{\mathrm{g}}$ bins, giving $\tilde{M}_d(k)$. The
estimates are then shrunk toward the next coarser level:
\begin{align}
  P_1(k) &= \frac{\tilde{M}_1(k) + \lambda\,\epsilon_0}{n_1(k)+\lambda},
  \label{eq:shrink1}\\
  P_d(k) &= \frac{\tilde{M}_d(k) + \lambda\,P_{d-1}(k')}{n_d(k)+\lambda},
  \qquad d=2,3,4,
  \label{eq:shrinkd}
\end{align}
where $k'$ is the parent key of $k$, $\lambda>0$ controls the strength of
shrinkage, $\epsilon_0$ is a small floor, and all entries are clipped to
$[0,1]$. For a hop with context $x$, the predictor uses the finest level whose key
was observed in training, $\hat{p}(c,b\mid x)=P_{d^\ast}(k_{d^\ast}(x))[c,b]$;
a source site never observed in training receives the flat value $\epsilon_0$.
Entries estimate the probability that a pair belongs to the visibility set, not
that the object is first seen there, and need not sum to one. For each key, the
entries are stored in sorted order, so the cost of any threshold follows from a
binary search. The calibration layer adds $O(|\mathcal{G}(x)|\,m)$ operations per
hop.

% verbatim from 7c55fbc sections/sec4_method.tex (definition of the uniform scale; verdict 11 §1-3)
\paragraph*{Uniform scale (Section IV of version 7c55fbc)}
The original \emph{uniform} scale corresponds to
\begin{equation}
  \tau_{\mathrm{uni}}(q) = \hat{F}^{-1}(1-q),
  \label{eq:tau-uni}
\end{equation}
so that $q$ is a coverage level on the calibration data.

\paragraph*{Why a logarithmic scale}
On the uniform scale, a bucket spans $1/m=\mvpBucketWidth$ in coverage level, so near the
target one bucket changes the miss level by up to a factor of two, and $q=1$
searches all candidates. Since cost grows faster than linearly as the miss level
falls, MVP's randomization over this coarse grid raises the average cost. On the
logarithmic scale, one bucket changes the miss level by the constant factor
$a_{\min}^{1/m}\approx\logBucketFactor$, which refines the grid where the target lies,
and the most permissive threshold is bounded at
$\tau_{\min}=\hat{F}^{-1}(a_{\min})$. the ablation in Section~\ref{app:heldout} of the supplementary material separates the
two effects. The bound has a cost: hops with $\sigma<\tau_{\min}$ are missed
whatever MVP does, so a group can reach its target only if
\begin{equation}
  \Pr\bigl[\sigma<\tau_{\min}\,\big|\,x\in G\bigr] < \alpha .
  \label{eq:feasible}
\end{equation}
This holds on the calibration data, where the left-hand side is about $a_{\min}$,
but can fail when the predictor is misspecified (Section~\ref{app:heldout} of the supplementary material).


\paragraph*{MVP procedure}
Algorithm~\ref{alg:mvp} states the procedure of Section~\ref{M-sec:method:mvp},
including the warm start and the handling of delayed, partial feedback. The
learning rate is $\eta=\sqrt{\ln(|\mathcal{G}|m)/(2K_\epsilon|\mathcal{G}|m)}$ with
$K_\epsilon=\sum_{n\ge 0}f(n)^{-2}$.

\begin{algorithm}[!htbp]
\caption{MVP threshold calibration with delayed, partial feedback}
\label{alg:mvp}
\begin{algorithmic}[1]
\Require calibration scores and contexts in time order; parameterization
  $\tau(\cdot)$ from \eqref{eq:tau-uni} or \eqref{M-eq:tau-log}; $m$, $r$, $\eta$
\State $V_{G,i}\gets 0$, $n_{G,i}\gets 0$ for all $G\in\mathcal{G}$, $i\in[m]$
\For{each calibration hop in time order} \Comment{warm start}
  \State choose $q$ as in lines~\ref{alg:choose1}--\ref{alg:choose2};
  update $(V,n)$ immediately with its outcome
\EndFor
\For{each test hop $t$ in order of $t^{\mathrm{exit}}_t$}
  \For{each feedback hop $j$ resolved by $t^{\mathrm{exit}}_t$}
    \State $i\gets$ bucket of $q_j$; for all $G\ni x_j$:
      $n_{G,i}\mathrel{+}=1$,
      $V_{G,i}\mathrel{+}=\mathbb{1}[\sigma_j\ge\tau(q_j)]-(1-\alpha)$
  \EndFor
  \State compute $C_i(x_t)$ for $i\in[m]$ by \eqref{M-eq:mvp-C} \label{alg:choose1}
  \If{$C_i>0$ for all $i$} $q_t\gets 0$
  \ElsIf{$C_i<0$ for all $i$} $q_t\gets 1$
  \Else
    \State $i^\ast\gets$ first index with $C_{i^\ast}C_{i^\ast+1}\le 0$;
      $p\gets|C_{i^\ast+1}|/(|C_{i^\ast}|+|C_{i^\ast+1}|)$
    \State $q_t\gets i^\ast/m-1/(rm)$ w.p.\ $p$, else $q_t\gets i^\ast/m$
      \label{alg:choose2}
  \EndIf
  \State search $\mathcal{S}_{\tau(q_t)}(x_t)$
  \If{hop $t$ provides feedback}
    \State schedule its outcome at its arrival time if hit, else at
      $t^{\mathrm{exit}}_t+T_{\max}$
  \EndIf
\EndFor
\end{algorithmic}
\end{algorithm}

Table~\ref{tab:parameters} lists all parameter values.

% Flush the appendix floats before the references.
\clearpage


